\documentclass[11pt]{article}
\usepackage[T1]{fontenc}
\usepackage{amsmath,amssymb,hyperref,geometry}
\geometry{margin=1in}
\title{Native Feedback and Interaction Extension\\Research Contract v0.1}
\author{}\date{16 November 2026}
\begin{document}\maketitle
\paragraph{Purpose and status.}
Investigate a conservative extension of the presented relative context/content construction that supplies genuinely new feedback or interaction capabilities while retaining native witnesses and directed comparisons. This is a research target, not a theorem or a claim that a universal completion exists. It does not redefine completion of the existing feedback-and-concurrency contract. General admissible feedback, branching graded interaction and infinite completion remain the broader ambition; their logical dependencies must be derived rather than bundled into one presumed construction.

\section{Baseline and structured inputs}
Begin with the independently specified context and content presentations, pure interface, native grade maps and changing-endpoint body squares, and coherent crossing underlying
\[
 M(A,B)=(e_A\downarrow B),\quad N(A,B)=(A\downarrow d_B),\quad H(A,B)=(e_A\downarrow d_B).
\]
The underlying hom categories alone are not the input. State which admission diagrams, joint effects, coupling choices, reindexings and coherence data each proposed extension consumes. Supplied choices remain explicit data unless a justified universal property determines them up to the stated notion of equivalence.

Two baselines must be compared. The original construction already supports finitary polynomial occurrences under its hypotheses. The newer local-joint construction derives native coefficient lifts and shared-effect comparisons for an independently graded one-hole content family. Earlier addressed branching uses an independence policy not interchangeable with an interfering coupling. Neither ``branching exists'' nor ``a recursive runtime exists'' is the missing theorem. The unresolved combination is native multi-occurrence substitution with chosen shared effects and its compatibility with feedback. Observable-cover and residual-profile certificates preserve a declared test interface; they do not construct this combination or replace its native data.

\section{First investigation: native branching with interference}
The initial route is a general construction or obstruction for branching graded occurrences over retained shared interfaces. Specify the allowed occurrence shapes, changing grades and interfaces, dependent continuations and joint admission. Derive native substitution and comparison transport from independently checkable local data, including their unit, associativity, reindexing and interaction coherence. Distinguish coherent constructors from pointwise existence of lifts. A whole-system comparison equation supplied as a hypothesis is not its derivation.

Investigate compatibility with retained-state closure and guarded recurrence already available. Do not duplicate a shared state into independent copies, replace chosen joint effects by a product law, or encode branching as a unary process and count that encoding alone as new categorical generality. A parameterized theorem may impose justified conditions, but these must expose the compatibility mechanism rather than merely select a convenient example. If evidence favors a different first route, record the mathematical reason and its relation to the broader ambition; this ordering is not the ultimate definition of the target.

\section{Extension interface and conservativity}
Before claiming a completion, specify its structured sources, targets and structure-preserving maps, with the comparison cells needed for its coherence. Distinguish a formally generated operation from an interpretation of that operation in a semantic class. For each feedback operation claimed, state its typed domain of admissible connections and how admissibility behaves under substitution and native comparison.

Declare precisely which old objects, arrows, equations and operations are preserved and reflected, whether new comparisons between old objects are permitted, and how distinct native maps invisible to a carrier interpretation remain distinguished. Prove the selected conservativity statement. Do not assume local fullness or a total trace merely from the word ``conservative.'' Any universal property is relative to the declared class of targets and maps; no strongest completion across unspecified doctrines is required.

\section{Independent continuation directions}
Retained-state closure, seeded delayed feedback, instantaneous cyclic interconnection and general iteration are different operations. Additional exactness, guardedness, continuity or fixed-point assumptions must be attached to the laws that use them; directed cells are not inverted without proof. The first branching result need not settle instantaneous feedback.

Infinite raw histories, compatible finite effect laws, divergence, and an infinite native coefficient object are likewise distinct. If a new operation needs infinite semantic objects, state the required existence and transport hypotheses and prove the completion or its obstruction. An already-known finite-support failure is baseline evidence, not a new result. Neither this direction nor arbitrary instantaneous feedback is a compulsory dependency of every useful finite native extension.

\section{Completion evidence for the selected investigation}
A positive outcome supplies a precise structured extension problem, a construction with proved coherence and the declared conservativity, and a substantive interpretation exhibiting a capability not supplied by the pinned baseline. Prefer a realization combining nontrivial native grades, branching recurrence and interfering effects, rather than disconnected examples. Explain whether the separating claim is nondefinability, failure of interpretation, or nonderivability from old hypotheses, and provide evidence at that strength. A free syntax without the needed interpretation or a change of presentation alone is insufficient.

An obstruction outcome supplies a counterexample at the proposed extension interface, identifies the failed obligation, and proves a meaningful repaired or restricted alternative where possible. State exactly which class is ruled out; failure of one candidate is not failure of all completions. A previously recorded obstruction alone does not discharge this investigation. If neither outcome is obtained, report the unresolved obligation without declaring completion.

Use existing constructions and primary literature where they supply the needed machinery; separate reuse, adaptation and potentially new claims. Obtain independent review of load-bearing written arguments. Keep formal proof, finite diagnostics and written proof distinct. No Lean installation, formalization campaign or TeX compilation is required for this investigation. Preserve source checkpoints and reviews outside the transient workspace.

\paragraph{Authoritative starting sources (paths relative to the contextual-composition project).}
\path{synthesis/Integrated_Relative_Computation_Acceptance_Dossier.tex};
\path{foundations/local-joint-comparison/Local_Joint_Comparison_and_Generated_Interference.tex};
\path{feedback-and-concurrency/observable-covers/Observable_Covers_and_Local_Observers.tex};
\path{feedback-and-concurrency/observable-covers/Residual_Profiles_and_Admission_Reflection.tex};
\path{contracts/Feedback_and_Concurrency_Research_Contract_v0_1_3.tex}.
Pin the versions used by each mathematical result; this contract supplies a target, not permanent assumptions about later source changes.
\end{document}
